\chapter{Code dans l'ESP32}
\label{sec:CodeESP32}
Le code de l'ESP32 est principalement employé pour la communication entre lui et le DSP et aussi pour host l'interface web. Cette section vient donc détailler les différents éléments de son code.

\subsection{Poids des codes}
À partir du compiler d'Arduino, il est possible de réaliser une vérification de l'espace prise dans la mémoire de l'ESP32. De cela on obtient les informations suivantes :

\begin{quote}
  Le croquis utilise 979768 octets (74\%) de l'espace de stockage de programmes. Le maximum est de 1310720 octets.

  Les variables globales utilisent 45324 octets (13\%) de mémoire dynamique, ce qui laisse 282356 octets pour les variables locales. Le maximum est de 327680 octets.
\end{quote}

D'un point de vue de la taille des fichiers, on a pour le code de l'ESP32 mesure 9.9 ko et celui du site web 10.6 ko.

\subsection{Code arduino}
Le code arduino est un petit code contenant de nombreuse fonction pour la communication ainsi que l'affichage de l'interface web.\\

Les librairies employé pour réaliser cela sont les suivantes :

\begin{minted}{c++}
#include <WiFi.h>
#include <WebServer.h>
#include <FastLED.h>
#include <HardwareSerial.h>
#include <stdlib.h>  // Pour atof()
#include <string.h>  // Pour strtok()

#include "WebPage.h"
\end{minted}

Elle contient les librairies principales suivantes :

\begin{itemize}
  \item WiFi.h : Une librairie de l'ESP32 permettant la gestion du WiFi. Elle sera utilisé pour transformer l'ESP32 en point d'accès.
  \item WebServer.h : Une librairie permettant d'heberger un site web et répondre à des requettes d'utilisateurs.
  \item HardwareSerial.h : Une librairie pour la communication en UART utilisé pour pour la communication entre l'ESP32 et le DSP.
  \item
\end{itemize}

Les principales définitions sont les suivantes :

\begin{minted}{c++}
// Nom et mot de passe du wifi crée par l'esp32
#define AP_SSID "ESP32-S2-WIFI"
#define AP_PASS "TFE-TB2025-HEIG"
// ... (LED setup) ...
HardwareSerial MySerial(1);
#define BAUDRATE 115200
#define TX1 17
#define RX1 18
// ... (Variables globales) ...
IPAddress Actual_IP;
IPAddress PageIP(192, 168, 4, 1);
// ...
WebServer server(80);
\end{minted}

Trois de ces définitions permettent de se connecter à l'interface web :

\begin{itemize}
  \item Nom du wifi : ESP32-S2-WIFI
  \item Mot de passe : TFE-TB2025-HEIG
  \item Lien du site web : 192.168.4.1
\end{itemize}

La fonction setup permet d'inialiser correctement l'interface web avec les importantes initialisaions suivantes :

\begin{minted}{c++}
void setup() {
  // ... (Initialisations Serial et LED) ...
  MySerial.begin(BAUDRATE, SERIAL_8N1, RX1, TX1, false);
  // ...
  WiFi.softAP(AP_SSID, AP_PASS);
  WiFi.softAPConfig(PageIP, gateway, subnet);
  // ...
  server.on("/", SendWebsite);
  server.on("/xml", SendXML);
  server.on("/BUTTON_START", ProcessButton_START);
  server.on("/BUTTON_STOP", ProcessButton_STOP);
  server.begin();
}
\end{minted}

La boucle principale permettant d'afficher le site web et de communiquer avec le DSP est la suivante :

\begin{minted}{c++}
void loop() {
  // ... (Réception de l'UART) ...
  while (MySerial.available() > 0) {
    uint8_t c = MySerial.read();
    if (RxIndex < UART_BUFFER) { buffer[RxIndex++] = c; }
    // ...
  }

  // Traitement des données
  if (RxIndex > 0) {
    buffer[RxIndex] = '\0';
    if (buffer[0] == 0x02 && buffer[RxIndex - 1] == 0x03) {
      char* token = strtok(&buffer[1], ",");
      int i = 0;
      while (token != NULL && i < 4) {
        floatValues[i++] = atof(token);
        token = strtok(NULL, ",");
      }
      updateValues();
      RxIndex = 0;
    }
    // ...
  }
  server.handleClient(); // Absolutely necessary
}
\end{minted}

Un bouton interfactif est actuellement mise en place comme suit :

\begin{minted}{c++}
void ProcessButton_START() {
  btn_state = 1;
  MySerial.print('\x02');
  MySerial.print(btn_state);
  MySerial.print('\x03');
  leds[0] = CRGB::Red;
  FastLED.show();
  server.send(250, "text/plain", ""); //Send web page
}
// (ProcessButton_STOP est similaire)
\end{minted}

L'envoie de donnée se réa

\subsection{Code de l'interface web}
L'interface se présente comme suit :



Le code de l'interface se compose en 3 parties :

\begin{itemize}
  \item Un code HTML pour
  \item Un code CSS pour
  \item Un code JavaScript pour
\end{itemize}

Ces 3 parties seront détaillé dans les sous-sections qui suivent.

\subsubsection{HTML}

\subsubsection{CSS}

\subsubsection{JavaScript}

\section{A faire}
\begin{itemize}
  \item Ajouter la mesure de courant
  \item Régler le chargement infinie de la page
  \item Créer un QR code pour se connecter à la page
  \item 
\end{itemize}
